% Propietario conservado: /Users/ruben/Documents/New project/output/EXTREMA_MEDIA_RAZON_RECIPROCIDAD_ES_EN_20260918/ARTICULO/ES/source/libro/reciprocidad_positiva.tex
\chapter{Positive deformation and lattice duality of the generated register}
\label{x_reciprocidad:rec:positiva}

\section{Coupling direction and unimodular flow}

Extraction of the register, its rational reading and the incidence projectors have been constructed in the preceding sections. Base one thousand, the phase origin, carry boundary and chart of the twelve positions are preserved. In that chart,
\[
 K=(234,543,140,729,659,824,621,058,914,794,146,601),
 \qquad u=P_3P_{11}K,
\]
\[
 P_{11}=I-\frac1{12}\mathbf1\mathbf1^{\mathsf T},\qquad
 u\perp\mathbf1,\qquad
 \|u\|^2=\frac{6638585+2275584\sqrt5}{20}>0.
\]
The finite definition of \(P_3\), the order of its classes and the exact evaluation of this norm appear in Lemma~\ref{x_reciprocidad:lem:k-sector-no-nulo}. The following realization uses that generated object. We write
\begin{equation}
 \rho=\frac{\pi_{\HMT}}{\varphi_{\HMT}^2}>0,\qquad
 h_i=(\log\rho)\frac{u_i}{\|u\|},\qquad
 r_i(t)=e^{th_i},\qquad g_t=\operatorname{diag}(r_i(t)).
 \label{x_reciprocidad:rec:eq:flujo}
\end{equation}
The regional coordinates are evaluated before this composition. The parameter \(t\in\RR\) indexes a mathematical deformation; its definition does not identify it with the TPK clock or with a refinement depth.

\begin{proposition}[Scale compensation]
\label{x_reciprocidad:rec:prop:flujo}
The flow satisfies
\[
 g_tg_s=g_{t+s},\qquad g_0=I,\qquad g_t^{-1}=g_{-t},\qquad
 \det g_t=\prod_i r_i(t)=1.
\]
\end{proposition}
\begin{proof}
Orthogonality to the uniform mode gives \(\sum_i h_i=0\). Multiplication of exponentials proves the group law and inversion. The determinant is \(\exp(t\sum_i h_i)=1\). Each diagonal factor is positive.
\end{proof}

Conservation is multiplicative and global: each radius may vary. In logarithmic coordinates the zero sum is preserved. This distinction will allow the transverse projector of rank ten to be interpreted geometrically.

\section{Exact reciprocity of the double circle}
\label{x_reciprocidad:rec:doble-circulo}

Consider the complex chart and its inverse
\[
 z(s)=\frac{s-1}{s},\qquad s(z)=\frac1{1-z}.
\]
For \(r>0\), \(z=re^{i\theta}\ne1\), define \(s(r,\theta)=(1-re^{i\theta})^{-1}\).

\begin{proposition}[Conjugation of the spectral mirror]
\label{x_reciprocidad:rec:prop:espejo}
The following hold
\begin{align}
 s(r^{-1},\theta)&=1-\overline{s(r,\theta)},\label{x_reciprocidad:rec:eq:espejo}\\
 \Re s(r,\theta)-\frac12&=
 \frac{1-r^2}{2|1-re^{i\theta}|^2}.\label{x_reciprocidad:rec:eq:lado}
\end{align}
The unit circle, except for the excluded point \(z=1\), corresponds to the line \(\Re s=1/2\).
\end{proposition}
\begin{proof}
Radial inversion at a fixed angle is \(z'=1/\overline z\). Therefore,
\[
 \frac1{1-z'}=\frac{\overline z}{\overline z-1}
 =1-\frac1{1-\overline z}.
\]
Moreover, \((1-z)^{-1}=(1-\overline z)/|1-z|^2\). Subtracting \(1/2\) from its real part produces \eqref{x_reciprocidad:rec:eq:lado}.
\end{proof}

The relation to the radii of \eqref{x_reciprocidad:rec:eq:flujo} is explicit: \(r_i(-t)=r_i(t)^{-1}\). Equality of products preserves volume, while the proposition identifies an involution between complex charts. Locating the zeros of a function requires the properties of that function; it is not a consequence of the unit product of radii.

Three circular realizations in the corpus should be distinguished. The spectral chart takes \(s=1/2+i\gamma\) to \(z_\gamma=(\gamma+i/2)/(\gamma-i/2)\). The dual circle of integer memory takes an integer \(n\) to the character \(z^n\). Nonadic compression takes \(z\) to \((8+z)/9\), a circle with center \(8/9\) and radius \(1/9\), tangent to the unit circle at \(1\). The modal ratio \(5/9\) belongs, by contrast, to the eigenvalues \(9,5\) of the APP block \(\left(\begin{smallmatrix}7&2\\2&7\end{smallmatrix}\right)\). Their functions remain distinct when composed.

\section{Positroid strata and balanced invariants}

Let \(X\) be a real matrix of rank \(k\), with twelve columns, and let \(\Delta_I(X)\) be its maximal minor with columns \(I\). The right action satisfies
\begin{equation}
 \Delta_I(Xg_t)=e^{t\sum_{i\in I}h_i}\Delta_I(X).
 \label{x_reciprocidad:rec:eq:plucker}
\end{equation}
Diagonal positivity preserves all signs and all vanishing minors. Thus, for \(1\le k\le11\), a trajectory starting in \(\operatorname{Gr}_{\ge0}(k,12)\) remains in the same positroid stratum. Stratification by minors is used here as a subsequent realization; see Postnikov's general construction~\cite{x_reciprocidad:rec:postnikov}.

\begin{proposition}[Positive duality of complementary ranks]
Let \(A_{\pm}=\operatorname{diag}(1,-1,1,-1,\ldots)\) and
\(\mathfrak D(V)=V^\perp A_{\pm}\), with the positive orientation of their complementary minors. Then
\[
 \mathfrak D(Vg_t)=\mathfrak D(V)g_{-t}.
\]
\end{proposition}
\begin{proof}
Since \(g_t\) is diagonal and self-adjoint,
\((Vg_t)^\perp=V^\perp g_t^{-1}\). The matrix \(A_{\pm}\) commutes with \(g_t\). For the sign of the minors, the orthogonal complement introduces, up to a global sign, \((-1)^{\sum_{j\in J}j}\). Multiplication of the columns by \((-1)^{j-1}\) leaves a sign independent of \(J\); it is removed by a single global orientation. The complementary minors are then nonnegative.
\end{proof}

If \(\mathbf1_I+\mathbf1_J=\mathbf1_L+\mathbf1_M\) and the denominator is nonzero, coordinatewise cancellation of the exponents proves
\begin{equation}
 \frac{\Delta_I(Xg_t)\Delta_J(Xg_t)}{\Delta_L(Xg_t)\Delta_M(Xg_t)}
 =\frac{\Delta_I(X)\Delta_J(X)}{\Delta_L(X)\Delta_M(X)}.
 \label{x_reciprocidad:rec:eq:cocientes}
\end{equation}
In rank two, the four-index quotients and their exchange relations are recovered. By contrast, the product of complementary minors of the two paired trajectories receives the same factor twice and may vary. Nor does the circulant frame \(O_K\), invertible for recovery, acquire the property of being totally nonnegative through this action: its minors preserve the signs established in its own realization.

\section{Covolume, dual lattice and ternary symmetry}

Let \(\Lambda\) be the previously constructed marked Leech lattice, self-dual and of covolume one in the metric of its twelve orthogonal planes of type \(A_2\). The marked isometry
\[
 \sigma=\bigoplus_{i=1}^{12}P_i\mathsf C P_i^{-1},\qquad
 \mathsf C^2+\mathsf C+I=0,
\]
preserves the gluing and \(\Lambda\), has order three and has no nonzero fixed vectors. The \(P_i\) are frames of the planes, distinct from the projectors \(P_3,P_{11}\). Define
\begin{equation}
 G_t=\bigoplus_{i=1}^{12}r_i(t)I_{A_{2,i}},\qquad
 \Lambda_t=G_t\Lambda.
 \label{x_reciprocidad:rec:eq:red}
\end{equation}

\begin{theorem}[Reciprocal lattice deformation]
\label{x_reciprocidad:rec:thm:red}
For every \(t\in\RR\),
\[
 \operatorname{covol}(\Lambda_t)=1,\qquad
 \Lambda_t^*=\Lambda_{-t},\qquad \sigma\Lambda_t=\Lambda_t.
\]
Every isometry of \(\Lambda\) commuting with \(G_t\) also preserves \(\Lambda_t\).
\end{theorem}
\begin{proof}
Each radius acts on two coordinates; thus,
\(\det G_t=\prod_i r_i(t)^2=1\). For an invertible operator \(G\), a vector \(y\) belongs to \((G\Lambda)^*\) exactly when \(G^*y\in\Lambda^*\). Therefore
\((G\Lambda)^*=G^{-*}\Lambda^*\). Applying self-duality, self-adjointness and inversion yields \(\Lambda_t^*=G_{-t}\Lambda\). Finally, \(G_t\) is scalar on each plane and commutes with \(\sigma\), whence
\(\sigma G_t\Lambda=G_t\sigma\Lambda=G_t\Lambda\). The same calculation applies to the other commuting isometries. The operator \(\sigma\) does not change, so it preserves its order and absence of fixed vectors.
\end{proof}

The point \(t=0\) recovers the original lattice with its integrality, parity, self-duality and exceptional realization. The complete family is a family of unimodular Euclidean lattices in the sense of covolume one; away from the original point, neither integrality nor a Moonshine realization is assumed. The theorem identifies exactly the ternary symmetry that does persist and the duality relating the two directions of deformation.

\begin{figure}[H]
\centering
\begin{tikzpicture}[node distance=3mm,>=Latex,every node/.style={font=\fontsize{8.5}{10}\selectfont,align=center},box/.style={draw=ink,rounded corners=2pt,text width=0.16\linewidth,inner sep=3pt,minimum height=25mm}]
\node[box] (a) {APP--TRIT--TPK: enriched state and joint continuum};
\node[box,right=of a] (b) {Register \(K\), regional publications and bilateral archive};
\node[box,right=of b] (c) {Recoverable direction \(u_K\) and positive flow of determinant one};
\node[box,right=of c] (d) {Conjugate radii, dual lattices and hyperbolic charts};
\node[box,right=of d] (e) {Spectral reciprocity and oriented recovery of action};
\draw[->] (a)--(b);\draw[->](b)--(c);\draw[->](c)--(d);\draw[->](d)--(e);
\end{tikzpicture}
\caption{Dependencies of the extension. The arrows represent compositions specified in the proofs; recovery preserves the charts and orientation.}
\label{x_reciprocidad:fig:rec-causal}
\end{figure}
